\documentclass{article}
\usepackage{amsmath,amssymb}
\usepackage[T1]{fontenc}
\usepackage[utf8]{inputenc}
\usepackage[english]{babel}
\usepackage{graphicx}
\usepackage{hyperref}
\usepackage{booktabs}

\title{Menu Intelligence: Structured Menu Extraction and Categorization}
\author{Philipp Philippov}
\date{\today}

\begin{document}
\maketitle

\begin{abstract}
This report describes Menu Intelligence, a web service for turning raw menu
text or menu documents into a structured menu representation. The project is
framed as a pipeline of NLP tasks: menu item categorization, structured field
extraction, and final menu aggregation. The repository is available at
\url{https://github.com/philippovdev/nlp-menu-intelligence}.

Main points to fill:
\begin{itemize}
  \item dataset size and source summary
  \item strongest baseline and final model
  \item key quantitative results
\end{itemize}
\end{abstract}

\section{Introduction}

Explain why menu structuring matters for restaurant tooling, digital menus, and
catalog maintenance. State that the project targets both a real deployable
service and a measurable NLP pipeline.

\subsection{Team}

State the team members and responsibilities. If the project is solo, describe
the roles covered by the single contributor.

\section{Related Work}

The closest published work to this project does not come from one single
benchmark. Instead, it splits across three literatures. First, short text and
product-title categorization papers motivate efficient text baselines and show
that catalog-style titles can be classified with strong simple models or
specialized neural architectures \cite{joulin-etal-2017-bag,xia-etal-2017-large,cevahir-murakami-2016-large}.
Second, food-domain NLP is real but fragmented: the nearest tasks are food-item
healthiness classification, dish name recognition from recipe text, and
large-scale menu normalization or clustering
\cite{wiegand-klakow-2013-towards,hu-etal-2023-diner,latif-etal-2025-restaurant}.
Third, structured extraction work is better represented by restaurant order
extraction and document or receipt understanding resources such as FUNSD, CORD,
and SROIE \cite{peng-etal-2017-may,jaume2019funsd,park2019cord,huang2019icdar2019}.

The important gap is that I did not find a public paper that matches the exact
task shape of this project: short menu item text, fixed menu taxonomy labels,
deterministic slot extraction, and later OCR-noisy document inputs. The
literature therefore supports the project mostly through adjacent tasks rather
than one established benchmark.

\begin{table}[t]
\centering
\begin{tabular}{p{0.18\linewidth}p{0.15\linewidth}p{0.16\linewidth}p{0.12\linewidth}p{0.14\linewidth}p{0.17\linewidth}}
\toprule
Approach & Task & Dataset & Metric & Result & Notes \\
\midrule
fastText \cite{joulin-etal-2017-bag} & short text classification & AG / DBPedia and related benchmarks & Accuracy & 92.5 on AG with bigrams & strong CPU baseline motivation \\
Xia et al. \cite{xia-etal-2017-large} & product title categorization & millions of Japanese titles, 35 classes & Accuracy & $>$96\% & close product-title analog \\
Hu et al. \cite{hu-etal-2023-diner} & dish name recognition & DiNeR & F1 / EM & 59.2 / 23.2 on OOD split & food-domain analog, not menu categorization \\
Latif et al. \cite{latif-etal-2025-restaurant} & menu clustering & $>$700k menu items & similarity / coverage & 0.98 intra-cluster sim.; 100\% coverage & closest menu-domain paper, but clustering not fixed-label classification \\
\bottomrule
\end{tabular}
\caption{Related work summary.}
\label{tab:related-work}
\end{table}

\section{Model Description}

Describe the working pipeline:
\begin{itemize}
  \item document ingestion
  \item text extraction
  \item line normalization
  \item category prediction
  \item structured field extraction
  \item JSON aggregation and human review
\end{itemize}

Include a pipeline diagram here when it is ready.

\section{Dataset}

Describe:
\begin{itemize}
  \item source collection policy
  \item annotation format
  \item split policy
  \item category labels
  \item extraction labels
\end{itemize}

The current training-ready item-level dataset is now available as
`items.v2.jsonl` with a paired source manifest and generated statistics
artifact. This subset contains 432 annotated item-level menu entries from 12
public restaurant or hotel dining sources, uses source-level split assignment
to reduce leakage between similar items, and keeps public source references in
the paired manifest. The released `v2` file is a source-grounded synthetic
expansion rather than a source-faithful raw crawl: item texts are normalized
into consistent English surface forms and prices are normalized to `RUB` for a
public-repo-safe classification dataset. The earlier `v1` subset remains useful
as the first small classical-baseline set, but `v2` is the fixed dataset
intended for the next transformer training step.

\begin{table}[t]
\centering
\begin{tabular}{lccc}
\toprule
Statistic & Train & Valid & Test \\
\midrule
Restaurants & 8 & 2 & 2 \\
Documents & 8 & 2 & 2 \\
Items & 288 & 72 & 72 \\
Avg. tokens per item & 8.70 & 8.72 & 8.68 \\
Categories & 12 & 12 & 12 \\
\bottomrule
\end{tabular}
\caption{Dataset statistics.}
\label{tab:dataset-stats}
\end{table}

\section{Experiments}

\subsection{Metrics}

Define the metrics used for category classification and extraction:
\begin{itemize}
  \item Accuracy
  \item Macro-F1
  \item extraction Precision, Recall, and F1
\end{itemize}

\subsection{Experiment Setup}

Describe:
\begin{itemize}
  \item train, valid, and test split
  \item hardware
  \item run configuration
  \item model selection policy
\end{itemize}

For the current fixed classification comparison, all trained models use the
source-level `items.v2.jsonl` split with 288 training items, 72 validation
items, and 72 held-out test items. The first transformer run uses
`distilbert-base-uncased` with a sequence-classification head, maximum sequence
length 64, batch size 16 for training, batch size 32 for evaluation, learning
rate $2 \times 10^{-5}$, weight decay 0.01, warmup ratio 0.1, and early
stopping on validation Macro-F1.

\subsection{Baselines}

List the baselines used in the comparison:
\begin{itemize}
  \item heuristic pipeline
  \item TF-IDF + Logistic Regression
  \item TF-IDF + Linear SVM
  \item DistilBERT sequence classifier
\end{itemize}

\section{Results}

The first fixed-split comparison on the training-ready `items.v2.jsonl` dataset
shows that simple TF-IDF baselines remain stronger than the first transformer
run. On the held-out test split, TF-IDF + Logistic Regression reaches 0.7222
accuracy and 0.7009 Macro-F1, while TF-IDF + Linear SVM reaches 0.7083
accuracy and 0.6930 Macro-F1. The first transformer run, a fine-tuned
DistilBERT classifier, reaches 0.5556 accuracy and 0.5620 Macro-F1 on the same
test split.

The heuristic pipeline is still useful as an end-to-end reference because it
also evaluates price and size extraction, but its current `items.v2` artifact
is measured on the full dataset rather than on the fixed train/valid/test
classification split. On that full item set it reaches 0.6181 category
accuracy, 0.6466 Macro-F1, and exact-match price and size accuracy of 1.0000.

On the validation split, TF-IDF + Linear SVM gives the highest score at 0.7500
accuracy and 0.7502 Macro-F1, slightly above TF-IDF + Logistic Regression at
0.7361 / 0.7302. DistilBERT validation performance is lower at 0.6389 / 0.6248,
so the first transformer result does not support replacing the classical
baselines yet.

\begin{table}[t]
\centering
\begin{tabular}{lcccc}
\toprule
Method & Accuracy & Macro-F1 & Extraction F1 & Notes \\
\midrule
Heuristic baseline & 0.6181 & 0.6466 & N/A & full items.v2 set; price exact = 1.0000, size exact = 1.0000 \\
TF-IDF + Logistic Regression & 0.7222 & 0.7009 & N/A & items.v2 test split; valid = 0.7361 / 0.7302 \\
TF-IDF + Linear SVM & 0.7083 & 0.6930 & N/A & items.v2 test split; valid = 0.7500 / 0.7502 \\
DistilBERT sequence classifier & 0.5556 & 0.5620 & N/A & items.v2 test split; valid = 0.6389 / 0.6248 \\
\bottomrule
\end{tabular}
\caption{Main experimental results.}
\label{tab:main-results}
\end{table}

Interpret:
\begin{itemize}
  \item which methods work best and why
  \item which categories remain difficult
  \item where OCR quality limits downstream NLP
\end{itemize}

\begin{table}[t]
\centering
\begin{tabular}{p{0.9\linewidth}}
\toprule
Sample output goes here. \\
\bottomrule
\end{tabular}
\caption{Output samples.}
\label{tab:output}
\end{table}

\section{Conclusion}

Summarize:
\begin{itemize}
  \item dataset created
  \item baselines evaluated
  \item best-performing approach
  \item remaining work and future improvements
\end{itemize}

\bibliographystyle{apalike}
\bibliography{lit}

\end{document}
